% GENERATED FILE. Edit canonical lessons, then run npm run generate:books.
% canonical-source-hash: fnv1a64:7d30d78578f49947
% canonical-lessons: IT-C414-statale, IT-C414-provinciale, IT-C414-curva, IT-C414-buca, IT-C414-sterrato

\chapter{Strade di campagna}
\label{ch:it-strade-di-campagna}
\hlchaptermodality{414}

\begin{chapteropening}
\textbf{By the end of this chapter:} \emph{I can describe country roads: the main road, the provincial road, bends, potholes and dirt tracks.}

Complete the last lesson of chapter 414: I can describe country roads: the main road, the provincial road, bends, potholes and dirt tracks.
\end{chapteropening}

\section[la statale]{la statale — the main road, the state highway}

\label{lesson:IT-C414-statale}



\begin{quote}
Before the new one: say the Italian for the diversion, then the Italian for to divert.
\end{quote}

\subsection*{You'll want to know: la statale}
\textbf{la statale} — \textquotedblleft{}the main road, the state highway\textquotedblright{}.

\begin{quote}
Prendiamo la statale, l'autostrada costa troppo.
\end{quote}

\begin{quote}
\emph{Let us take the main road, the motorway costs too much.}
\end{quote}

Short for \emph{la strada statale}, a road run by the state.

\begin{cousinweb}
\emph{statale} is built on \emph{stato} 'state', a learned borrowing from Latin \emph{status} 'standing, condition', from \emph{stāre} 'to stand' (Lewis \& Short).
\end{cousinweb}

\begin{grammarlens}[title={What you've built}]
One more word for telling, explaining and describing at B1.
\end{grammarlens}

\subsection*{Your turn}
\begin{itemize}
\raggedright
  \item \emph{Say it:} \emph{la statale}
  \item \emph{Say it:} \emph{la statale}, once more
  \item \emph{Say it:} a sentence of your own with \emph{la statale}
  \item \emph{Recall:} say the Italian for the diversion, then the Italian for to divert, then say \emph{la statale} again
\end{itemize}

\begin{tcolorbox}[breakable,colback=teal!4,colframe=teal!35!black,title={Before you move on}]
What does la statale mean? (\textquotedblleft{}The main road, the state highway\textquotedblright{}.) Say it once more.
\end{tcolorbox}
\section[la provinciale]{la provinciale — the provincial road, the secondary road}

\label{lesson:IT-C414-provinciale}



\begin{quote}
Before the new one: say the Italian for to divert, then the Italian for the main road.
\end{quote}

\subsection*{You'll want to know: la provinciale}
\textbf{la provinciale} — \textquotedblleft{}the provincial road, the secondary road\textquotedblright{}.

\begin{quote}
La provinciale passa per tre paesi.
\end{quote}

\begin{quote}
\emph{The provincial road goes through three villages.}
\end{quote}

Short for \emph{la strada provinciale}; smaller than \emph{la statale}. As an adjective it can also mean narrow-minded.

\begin{cousinweb}
\emph{provinciale} is built on \emph{provincia} 'province', a learned borrowing from Latin \emph{prōvincia} 'a charge, a province' (Lewis \& Short).
\end{cousinweb}

\begin{grammarlens}[title={What you've built}]
One more word for telling, explaining and describing at B1.
\end{grammarlens}

\subsection*{Your turn}
\begin{itemize}
\raggedright
  \item \emph{Say it:} \emph{la provinciale}
  \item \emph{Say it:} \emph{la provinciale}, once more
  \item \emph{Say it:} a sentence of your own with \emph{la provinciale}
  \item \emph{Recall:} say the Italian for to divert, then the Italian for the main road, then say \emph{la provinciale} again
\end{itemize}

\begin{tcolorbox}[breakable,colback=teal!4,colframe=teal!35!black,title={Before you move on}]
What does la provinciale mean? (\textquotedblleft{}The provincial road, the secondary road\textquotedblright{}.) Say it once more.
\end{tcolorbox}
\section[la curva]{la curva — the bend, the curve}

\label{lesson:IT-C414-curva}



\begin{quote}
Before the new one: say the Italian for the main road, then the Italian for the provincial road.
\end{quote}

\subsection*{You'll want to know: la curva}
\textbf{la curva} — \textquotedblleft{}the bend, the curve\textquotedblright{}.

\begin{quote}
Rallenta, c'è una curva pericolosa.
\end{quote}

\begin{quote}
\emph{Slow down, there is a dangerous bend.}
\end{quote}

In a stadium, \emph{la curva} is the end stand where the keenest fans sit.

\begin{cousinweb}
\emph{curva} is inherited from Latin \emph{curva} 'curve' (EtymDB), from \emph{curvus} 'bent'.
\end{cousinweb}

\begin{grammarlens}[title={What you've built}]
One more word for telling, explaining and describing at B1.
\end{grammarlens}

\subsection*{Your turn}
\begin{itemize}
\raggedright
  \item \emph{Say it:} \emph{la curva}
  \item \emph{Say it:} \emph{la curva}, once more
  \item \emph{Say it:} a sentence of your own with \emph{la curva}
  \item \emph{Recall:} say the Italian for the main road, then the Italian for the provincial road, then say \emph{la curva} again
\end{itemize}

\begin{tcolorbox}[breakable,colback=teal!4,colframe=teal!35!black,title={Before you move on}]
What does la curva mean? (\textquotedblleft{}The bend, the curve\textquotedblright{}.) Say it once more.
\end{tcolorbox}
\section[la buca]{la buca — the pothole, the hole (in the ground)}

\label{lesson:IT-C414-buca}



\begin{quote}
Before the new one: say the Italian for the provincial road, then the Italian for the bend.
\end{quote}

\subsection*{You'll want to know: la buca}
\textbf{la buca} — \textquotedblleft{}the pothole, the hole (in the ground)\textquotedblright{}.

\begin{quote}
La strada è piena di buche.
\end{quote}

\begin{quote}
\emph{The road is full of potholes.}
\end{quote}

\emph{La buca} is a hole in the ground; \emph{il buco} is a hole in anything. \emph{La buca delle lettere} — the letterbox.

\begin{cousinweb}
\emph{buca} goes with \emph{buco} 'hole'. The origin is not settled; some link it to Latin \emph{bucca} 'cheek, mouth'.
\end{cousinweb}

\begin{grammarlens}[title={What you've built}]
One more word for telling, explaining and describing at B1.
\end{grammarlens}

\subsection*{Your turn}
\begin{itemize}
\raggedright
  \item \emph{Say it:} \emph{la buca}
  \item \emph{Say it:} \emph{la buca}, once more
  \item \emph{Say it:} a sentence of your own with \emph{la buca}
  \item \emph{Recall:} say the Italian for the provincial road, then the Italian for the bend, then say \emph{la buca} again
\end{itemize}

\begin{tcolorbox}[breakable,colback=teal!4,colframe=teal!35!black,title={Before you move on}]
What does la buca mean? (\textquotedblleft{}The pothole, the hole (in the ground)\textquotedblright{}.) Say it once more.
\end{tcolorbox}
\section[lo sterrato]{lo sterrato — the dirt road, the unpaved track}

\label{lesson:IT-C414-sterrato}



\begin{quote}
Before the new one: say the Italian for the bend, then the Italian for the pothole.
\end{quote}

\subsection*{You'll want to know: lo sterrato}
\textbf{lo sterrato} — \textquotedblleft{}the dirt road, the unpaved track\textquotedblright{}.

\begin{quote}
Dopo il paese c'è un lungo sterrato.
\end{quote}

\begin{quote}
\emph{After the village there is a long dirt road.}
\end{quote}

As an adjective: \emph{una strada sterrata}.

\begin{cousinweb}
\emph{sterrato} is built from \emph{s-} and \emph{terra} 'earth', which is inherited from Latin \emph{terra} (EtymDB): a road of bare earth.
\end{cousinweb}

\begin{grammarlens}[title={What you've built}]
One more word for telling, explaining and describing at B1.
\end{grammarlens}

\subsection*{Your turn}
\begin{itemize}
\raggedright
  \item \emph{Say it:} \emph{lo sterrato}
  \item \emph{Say it:} \emph{lo sterrato}, once more
  \item \emph{Say it:} a sentence of your own with \emph{lo sterrato}
  \item \emph{Recall:} say the Italian for the bend, then the Italian for the pothole, then say \emph{lo sterrato} again
\end{itemize}

\begin{tcolorbox}[breakable,colback=teal!4,colframe=teal!35!black,title={Before you move on}]
What does lo sterrato mean? (\textquotedblleft{}The dirt road, the unpaved track\textquotedblright{}.) Say it once more.
\end{tcolorbox}
